% FORMULARIO CONTROLLO AUTOMATICI
\documentclass[a4paper,11pt]{article}
\usepackage{amsmath,amssymb,siunitx,tikz}
\usepackage{geometry}
\usepackage{amsthm} % definisce ambienti theorem-like
\newtheorem{exercise}{Esercizio}[section]
\newtheorem{example}{Esempio}[section]

\geometry{margin=2cm}
\title{Formulario di Controlli Automatici}
\author{}
\date{\today}
\begin{document}
\maketitle
\tableofcontents
\newpage

%-----------------------------------------------------------------
\section*{Legenda}
\begin{description}
  \item[$s$] variabile complessa di Laplace
  \item[$G(s)$] funzione di trasferimento del plant
  \item[$H(s)$] funzione di trasferimento del sensore (feedback)
  \item[$K$] guadagno di retroazione
\end{description}

%=================================================================
\section{Modulo 1 – Fondamenti di sistemi}

% --------------------- Contenuti Modulo 1 ---------------------
\subsection*{Micro-lezione (15 righe)}
\begin{itemize}
  \item Un \emph{sistema} elabora un ingresso $u(t)$ e restituisce un'uscita $y(t)$.
  \item Se è \textbf{lineare} e \textbf{tempo invariante} (LTI) vale la sovrapposizione e poli/zeri non cambiano nel tempo.
  \item La \emph{quiete iniziale} impone condizioni iniziali nulle per studiare la sola risposta all’ingresso.
  \item Lo \textbf{schema a blocchi} rappresenta catene di funzioni di trasferimento e si riduce con 8 regole (somma, serie, parallelo, retroazione …).
  \item Nel \textbf{controllo in retroazione} si sottrae la misura dal riferimento: $e=r-y$.
  \item Il \textbf{controllo diretto} (feed–forward) agisce su $u$ senza chiudere l’anello ma richiede un modello preciso.
  \item \textbf{Stabilità interna}: tutti i poli di $G(s)$ nel semipiano sinistro $\Rightarrow$ risposta limitata.
  \item \textbf{Sensibilità} $S=\dfrac{1}{1+G(s)H(s)}$ misura l’effetto dei disturbi sull’uscita.
\end{itemize}

\subsection*{Esempio numerico – Cruise‑control}
Massa $m=1000\,\mathrm{kg}$, attrito lineare $b=50\,\mathrm{N\,s/m}$.
\[
G(s)=\frac{1/m}{s+\tfrac{b}{m}}=\frac{0{,}001}{s+0{,}05}.
\]
Tempo di assestamento ($\pm2\%$) I‑ordine: $t_s\approx4\,\tau=80\,\mathrm{s}$ con $\tau=20\,\mathrm{s}$.

\subsection*{Mappa concettuale}
\begin{center}\small
\begin{verbatim}
   r ─┬─►(+)──► C(s) ─► u ─► G(s) ─► y
      │          ^               │
      │          │               │
      └──────────┴── H(s) ◄──────┘
\end{verbatim}
\end{center}

\subsection*{Quiz rapido Mod. 1}
\begin{exercise}[V/F]
Indicare se le affermazioni sono vere (V) o false (F):
\begin{enumerate}
  \item Un sistema con polo in $s=0$ è asintoticamente stabile. \hfill\textbf{F}
  \item La somma di due sistemi LTI è ancora LTI. \hfill\textbf{V}
  \item In retroazione unitaria $T(s)=\dfrac{C(s)G(s)}{1+C(s)G(s)}$. \hfill\textbf{V}
\end{enumerate}
\end{exercise}

%-----------------------------------------------------------------
\section{Modulo 2 – Modelli dinamici \& Laplace}
\subsection*{Trasformate elementari (promemoria)}
\[
\mathcal{L}\{1\} = \frac{1}{s},\qquad
\mathcal{L}\{t\} = \frac{1}{s^{2}},\qquad
\mathcal{L}\{e^{at}\} = \frac{1}{s-a},\qquad
\mathcal{L}\{\sin \omega t\} = \frac{\omega}{s^{2}+\omega^{2}},\qquad
\mathcal{L}\{\cos \omega t\} = \frac{s}{s^{2}+\omega^{2}}.
\]

\subsection*{Antitrasformate rapide via scomposizione}
Dato
\[
Y(s)=\frac{k_1}{(s-a)}+\frac{k_2}{(s-b)}
\quad\Longrightarrow\quad
y(t)=k_1 e^{a t}+k_2 e^{b t}.
\]

\subsection*{Micro-lezione}
\begin{enumerate}
  \item Le equazioni differenziali lineari diventano algebriche applicando la trasformata di Laplace.
  \item Funzione di trasferimento: $G(s)=\dfrac{Y(s)}{U(s)}$ con condizioni iniziali nulle.
  \item Teorema del valore iniziale: $\displaystyle \lim_{t\to0^+}y(t)=\lim_{s\to\infty}sY(s)$.
  \item Teorema del valore finale (poli nel semipiano sinistro): $\displaystyle \lim_{t\to\infty}y(t)=\lim_{s\to0}sY(s)$.
  \item Convoluzione nel tempo $\Longleftrightarrow$ prodotto in dominio $s$.
\end{enumerate}

\subsection*{Esempio svolto}
Calcolare $y(t)$ per ingresso gradino unitario e
\[
G(s)=\frac{5}{s(s+2)}.
\]
\textbf{Soluzione}\\
\[
Y(s)=G(s)\frac{1}{s}= \frac{5}{s^{2}(s+2)}
=\frac{A}{s}+\frac{B}{s^{2}}+\frac{C}{s+2}.
\]
Con i residui si ottiene $A=\tfrac{5}{2}$, $B=-\tfrac{5}{2}$, $C=-\tfrac{5}{2}$:
\[
y(t)=\frac{5}{2}-\frac{5}{2}e^{-2t}-\frac{5}{2}t.
\]

\subsection*{Mini‑mappa}
\begin{center}\small
\begin{verbatim}
equaz.diff  →  Laplace  →  algebra  →  antitrasformata →  y(t)
\end{verbatim}
\end{center}

%-----------------------------------------------------------------
\section{Modulo 3 – Risposta al gradino}
% TODO

%-----------------------------------------------------------------
\section{Modulo 4 – Analisi armonica}
% TODO

%-----------------------------------------------------------------
\section{Modulo 5 – Stabilità \& retroazione}

\subsection*{Criterio di Routh – intervallo di stabilità (grado 3)}
Sia
\[
\chi(s)=a s^{3}+b s^{2}+c s + d = 0 ,\qquad a>0 .
\]
La \emph{tabella di Routh} è
\[
\begin{array}{c|cc}
s^{3} & a & c \\
s^{2} & b & d \\
s^{1} & x & 0 \\
s^{0} & d & 0
\end{array}
\qquad
x=\frac{b c - a d}{b}.
\]
Stabilità $\;\Longleftrightarrow\;$ tutti gli elementi della prima colonna hanno lo stesso segno, da cui
\[
a>0,\; b>0,\; d>0,\; x>0.
\]

\begin{example}[Intervallo di stabilità con parametro $K$]
Determinare i valori di $K$ che rendono stabile il sistema a retroazione unitaria con
\[
G_d(s)=K\,\frac{10^{4}}{s(s+5)(s+20)}.
\]
\textbf{Soluzione} – Equazione caratteristica:
\[
1+K\,\frac{10^{4}}{s(s+5)(s+20)}=0
\Longrightarrow
s^{3}+25 s^{2}+100 s + 10^{4} K = 0.
\]
Imponendo la positività della prima colonna della tabella di Routh si ottiene
\[
0 < K < 0{,}25.
\]
\end{example}

\begin{exercise}[Secondo esempio]
Ripetere la procedura per
\[
G_d(s)=K\,\frac{100}{s(s+3)(s+40)}\quad\Rightarrow\quad 0<K<51{,}6.
\]
\end{exercise}

\bigskip
\noindent\textbf{Formule utili (grado 3):}
\[
x=\frac{b c - a d}{b},\quad
y=\frac{b e - a f}{b},\quad
z=\frac{x d - b f}{x}.
\]
% TODO: verificare notazione ed esempi aggiuntivi

%-----------------------------------------------------------------
\section{Modulo 6 – Luogo delle radici}
% TODO

%-----------------------------------------------------------------
\section{Modulo 7 – Sintesi \& compensazione}
% TODO

\end{document}